\clearpage\begingroup
\hypertarget{native-328b}{}\HistoricalHead{323}{R}{BEOBACHTUNGEN AM REICHENBACHSCHEN MERIDIANKREISE.}\HistoricalFoot{}{41\textsuperscript{*}}
\begin{center}\textsc{Gauss} an \textsc{Bessel}, 28. Juni 1820.\end{center}
\noindent\hspace*{\parindent}. . . . . . Sehr interessant waren mir Ihre Beobachtungen von 6 Sternen. Ich schicke Ihnen hier meine Beobachtungen derselben Sterne. Es scheint, dass die meisten nahe mit den Ihrigen harmoniren und dass meine Beob[achtungen] die Polardistanzen nur etwa \(1''\) grösser geben mit Ausnahme des Sirius, wo meine Beob[achtungen] die Polardistanz wohl \(7''\) grösser geben.
\clearpage\SourceLeaf{329}\HistoricalHead{324}{L}{PRAKTISCHE UND SPHÄRISCHE ASTRONOMIE. BRIEFWECHSEL.}
\noindent Hier meine Beobachtungen[\begingroup\renewcommand{\thefootnote}{*)}\footnotemark\renewcommand{\thefootnote}{[*)}\footnotetext{Scheinbare Örter.]}\endgroup]. Das Datum ist durch die gezählten Sterntage designirt, so dass z. B. 53 die Culmination am 23. Februar bedeutet.
\begin{center}\begingroup\fontsize{9}{11}\selectfont\setlength{\tabcolsep}{2pt}
\begin{tabular}{@{}c|c|c|c|}\hline\hline
Capella&\(\beta\) Tauri&Sirius&Castor praec.\\\hline\hline
\begin{tabular}[t]{@{}rr@{}}
\multicolumn{2}{c}{I.}\\
53&\(5^\circ41'\ 4\overset{\prime\prime}{,}94\)\\
54&\(5{,}09\)\\
57&\(4{,}71\)\\
60&\(6{,}53\)\\
69&\(5{,}63\)\\[.8ex]
\multicolumn{2}{c}{II.}\\
79&\(5^\circ45'38\overset{\prime\prime}{,}35\)\\
85&\(38{,}62\)\\
88&\(41{,}23\)\\
95&\(41{,}75\)\\
100&\(39{,}94\)\\[.8ex]
\multicolumn{2}{c}{III.}\\
104&\(5^\circ41'\ 6\overset{\prime\prime}{,}52\)\\
108&\(6{,}43\)
\end{tabular}&
\begin{tabular}[t]{@{}rr@{}}
\multicolumn{2}{c}{I.}\\
71&\(23^\circ\ 2'38\overset{\prime\prime}{,}20\)\\
77&\(39{,}16\)\\
78&\(39{,}21\)\\[.8ex]
\multicolumn{2}{c}{II.}\\
85&\(23^\circ\ 7'11\overset{\prime\prime}{,}13\)\\
88&\(13{,}55\)\\
90&\(14{,}68\)\\
95&\(13{,}70\)\\[.8ex]
\multicolumn{2}{c}{III.}\\
106&\(23^\circ\ 2'39\overset{\prime\prime}{,}24\)\\
108&\(39{,}13\)\\
111&\(39{,}81\)\\
128&\(39{,}38\)
\end{tabular}&
\begin{tabular}[t]{@{}rr@{}}
\multicolumn{2}{c}{I.}\\
53&\(67^\circ58'10\overset{\prime\prime}{,}26\)\\
57&\(10{,}22\)\\
58&\(10{,}95\)\\
74&\(11{,}90\)\\
78&\(12{,}39\)\\[.8ex]
\multicolumn{2}{c}{II.}\\
81&\(68^\circ\ 2'45\overset{\prime\prime}{,}77\)\\
85&\(46{,}25\)\\
87&\(47{,}86\)\\
88&\(46{,}96\)\\
89&\(47{,}52\)\\
90&\(47{,}45\)\\
94&\(47{,}78\)\\
98&\(45{,}83\)\\
100&\(45{,}13\)\\
102&\(45{,}45\)\\[.8ex]
\multicolumn{2}{c}{III.}\\
104&\(67^\circ58'12\overset{\prime\prime}{,}09\)\\
106&\(10{,}81\)\\
126&\(8{,}08\)\\
133&\(6{,}00\)\\[.8ex]
\multicolumn{2}{c}{IV.}\\
150&\(68^\circ\ 2'39\overset{\prime\prime}{,}83\)
\end{tabular}&
\begin{tabular}[t]{@{}rr@{}}
\multicolumn{2}{c}{I.}\\
52&\(19^\circ13'\ 3\overset{\prime\prime}{,}28\)\\[.8ex]
\multicolumn{2}{c}{II.}\\
85&\(19^\circ17'33\overset{\prime\prime}{,}40\)\\
88&\(34{,}76\)\\
89&\(34{,}26\)\\
90&\(34{,}73\)\\
101&\(33{,}73\)\\[.8ex]
\multicolumn{2}{c}{III.}\\
106&\(19^\circ12'56\overset{\prime\prime}{,}97\)\\
107&\(56{,}65\)\\
112&\(59{,}00\)\\
115&\(58{,}80\)\\
130&\(57{,}30\)\\
137&\(58{,}65\)
\end{tabular}
\end{tabular}\par\endgroup\end{center}
\clearpage
\begin{center}\begingroup\fontsize{9}{11}\selectfont\setlength{\tabcolsep}{3pt}
\begin{tabular}{@{}c|c|c@{}}\hline\hline
Procyon&Pollux&Spica\\\hline\hline
\begin{tabular}[t]{@{}rr@{}}
\multicolumn{2}{c}{I.}\\
52&\(45^\circ48'49\overset{\prime\prime}{,}68\)\\
58&\(49{,}23\)\\[.8ex]
\multicolumn{2}{c}{II.}\\
85&\(45^\circ53'22\overset{\prime\prime}{,}00\)\\
88&\(23{,}94\)\\
89&\(22{,}71\)\\
90&\(23{,}58\)\\[.8ex]
\multicolumn{2}{c}{III.}\\
106&\(45^\circ48'46\overset{\prime\prime}{,}61\)\\
107&\(46{,}32\)\\
108&\(47{,}42\)\\
109&\(47{,}93\)\\
117&\(47{,}19\)\\
130&\(46{,}59\)\\
137&\(45{,}68\)\\
144&\(45{,}02\)\\[.8ex]
\multicolumn{2}{c}{IV.}\\
157&\(45^\circ53'19\overset{\prime\prime}{,}49\)\\
158&\(19{,}57\)
\end{tabular}&
\begin{tabular}[t]{@{}rr@{}}
\multicolumn{2}{c}{II.}\\
94&\(23^\circ\ 6'52\overset{\prime\prime}{,}82\)\\
95&\(51{,}55\)\\
101&\(51{,}73\)\\
102&\(51{,}08\)\\[.8ex]
\multicolumn{2}{c}{III.}\\
103&\(23^\circ\ 2'15\overset{\prime\prime}{,}47\)\\
104&\(16{,}24\)\\
111&\(15{,}18\)\\
115&\(15{,}20\)\\
116&\(15{,}08\)\\
126&\(15{,}95\)\\
130&\(14{,}35\)\\
137&\(15{,}51\)\\
144&\(15{,}31\)\\
145&\(15{,}40\)
\end{tabular}&
\begin{tabular}[t]{@{}rr@{}}
\multicolumn{2}{c}{III.}\\
108&\(61^\circ42'52\overset{\prime\prime}{,}27\)\\
109&\(52{,}70\)\\
116&\(52{,}77\)\\
132&\(54{,}14\)\\
143&\(55{,}36\)\\
144&\(53{,}70\)\\[.8ex]
\multicolumn{2}{c}{IV.}\\
162&\(61^\circ47'27\overset{\prime\prime}{,}39\)\\
169&\(25{,}84\)\\
177&\(26{,}25\)\\
178&\(25{,}40\)
\end{tabular}
\end{tabular}\par\endgroup\end{center}

\noindent\hspace*{\parindent}Ich habe auch die bisherigen Beob[achtungen] von \(\alpha\) Virginis beigefügt in der Hoffnung, dass dieses zur Aufklärung des grossen Unterschiedes beim Sirius etwas beitragen wird. Die Refractionen sind alle nach Ihrer Tafel (die ich mir in eine etwas andere Form gebracht habe) berechnet.

In den geraden Perioden habe ich gleich die Complemente zu \(360^\circ\) von dem, was das Instrument gab, angesetzt. Die scheinb[are] Zenithdistanz des Pols ist
\begin{center}\begingroup\setlength{\tabcolsep}{3pt}
\begin{tabular}{@{}l@{\quad}p{.57\textwidth}@{}}
\(\begin{array}{rrrr}\mathrm{I}&51^\circ&29'&31\overset{\prime\prime}{,}92\end{array}\)&\\
\(\left.\begin{array}{rrrr}\mathrm{II}&&34&6{,}46\\\mathrm{III}&&29&30{,}78\\\mathrm{IV}&&34&6{,}65\end{array}\right\}\)&sind noch keine Definitivresultate, werden aber schwerlich mehr als ein Paar Zehntel anders ausfallen\begingroup\renewcommand{\thefootnote}{**)}\footnotemark\endgroup.
\end{tabular}\par\endgroup\end{center}
\begingroup\renewcommand{\thefootnote}{**)}\footnotetext{Ausser der letzten (IV), die sich nur auf sehr wenige Circumpol[ar]sternbeob[achtungen] gründet und vielleicht etwa \(1''\) zu klein ist.}\endgroup
\clearpage\SourceLeaf{330}\HistoricalHead{325}{R}{BEOBACHTUNGEN AM REICHENBACHSCHEN MERIDIANKREISE.}
\noindent\hspace*{\parindent}Alle obigen Beob[achtungen], die der Capella ausgenommen, sind nach \textsc{Harding}s Hilfstafel berechnet. Ich setze bloss die Mittel her\begingroup\renewcommand{\thefootnote}{*)}\footnotemark\footnotetext{\textls[120]{Ohne eigene Bewegung} auf den Anf[ang] v[on] 1820 reducirt.}\endgroup:
\begin{center}\begingroup\fontsize{9}{11}\selectfont\setlength{\tabcolsep}{2pt}
\begin{tabular}{@{}c|c|c|}\hline\hline
\(\beta\) Tauri&Sirius\begingroup\renewcommand{\thefootnote}{**)}\footnotemark\endgroup&Castoris praec.\\\hline\hline
\begin{tabular}[t]{@{}lrr@{}}
I.&3&\(28^\circ26'41\overset{\prime\prime}{,}24\)\\
II.&4&\(41{,}75\)\\
III.&4&\(40{,}97\)\\[.8ex]
G[\textsc{auss}]&11&\(28^\circ26'41\overset{\prime\prime}{,}33\)\\
B[\textsc{essel}]&&\(42{,}55\)\\
Diff[erenz]&&\(+1{,}22\)
\end{tabular}&
\begin{tabular}[t]{@{}lrr@{}}
I.&5&\(-16^\circ28'35\overset{\prime\prime}{,}91\)\\
II.&10&\(36{,}01\)\\
III.&4&\(35{,}93\)\\
IV.&1&\(34{,}71\)\\[.8ex]
G.&20&\(-16^\circ28'35\overset{\prime\prime}{,}90\)\\
B.&&\(28{,}90\)\\
Diff.&&\(+7{,}00\)
\end{tabular}&
\begin{tabular}[t]{@{}lrr@{}}
I.&1&\(32^\circ16'20\overset{\prime\prime}{,}50\)\\
II.&5&\(21{,}67\)\\
III.&6&\(22{,}28\)\\[.8ex]
G.&12&\(32^\circ16'21\overset{\prime\prime}{,}88\)\\
B.&&\(23{,}20\)\\
Diff.&&\(+1{,}32\)
\end{tabular}
\end{tabular}\par\endgroup\end{center}
\begingroup\renewcommand{\thefootnote}{**)}\footnotetext{Aus meinen Beob[achtungen] des Sirius habe ich den wahrscheinl[ichen] Fehler Einer Beob[achtung] \(=0\overset{\prime\prime}{,}61\) gefunden.}\endgroup
\begin{center}\begingroup\fontsize{9}{11}\selectfont\setlength{\tabcolsep}{2pt}
\begin{tabular}{@{}c|c|c@{}}\hline\hline
Procyon&Pollux&Spica\\\hline\hline
\begin{tabular}[t]{@{}lrr@{}}
I.&2&\(5^\circ40'41\overset{\prime\prime}{,}07\)\\
II.&4&\(42{,}45\)\\
III.&8&\(42{,}39\)\\
IV.&2&\(43{,}39\)\\[.8ex]
G.&16&\(5^\circ40'42\overset{\prime\prime}{,}36\)\\
B.&&\(43{,}00\)\\
Diff.&&\(+0{,}64\)
\end{tabular}&
\begin{tabular}[t]{@{}lrr@{}}
II.&4&\(28^\circ27'5\overset{\prime\prime}{,}80\)\\
III.&10&\(6{,}41\)\\[.8ex]
G.&14&\(28^\circ27'6\overset{\prime\prime}{,}24\)\\
B.&&\(7{,}10\)\\
Diff.&&\(+0{,}86\)
\end{tabular}&
\begin{tabular}[t]{@{}lrr@{}}
III.&6&\(-10^\circ13'6\overset{\prime\prime}{,}67\)\\
IV.&4&\(4{,}22\)\\[.8ex]
G.&10&\(-10^\circ13'5\overset{\prime\prime}{,}69\)
\end{tabular}
\end{tabular}\par\endgroup\end{center}

\noindent\hspace*{\parindent}Im ganzen Junius ist das Wetter äusserst schlecht gewesen; nur ein Paar Circumpolarsternbeobachtungen habe ich erhalten. Ueberhaupt muss Ihr Klima \textls[120]{viel} besser sein, als das hiesige. Wenn ich den Ort des Pols bloss auf den Nordstern gründen wollte, so würde ich in Monatsfrist auch nicht eine einzige Bestimmung gehabt haben (kaum einmal eine \textls[120]{einseitige}). Für mich ist eine Anzahl von Normalsternen ein wahres Bedürfniss. Bei günstigem Wetter habe ich es übrigens recht gut thunlich gefunden, den Nordstern auch bei Tage aus einem Wasserspiegel zu beobachten; eine vollständige Beob[achtung] der Art, die ich am 13. Mai gemacht habe, werden Sie in den hiesigen G[elehrten] A[nzeigen, Werke VI, S. 429] gefunden haben. — Sonnenbeobachtungen am Kreise habe ich seit Eröffnung des Tagebuchs erst 29 erhalten. Die gleichzeitigen Beob[achtungen] am M[ittags]f[ernrohr] sind grösstentheils noch nicht reducirt (weil so häufig bei der Sonne oder den verglichenen Sternen ein oder der andere Faden ausfällt, so habe ich die Reduction aufgeschoben, bis ich die Fadenintervalle des neuen Netzes aus 25 vollständigen Nordst-[beobachtungen] haben würde, welche Anzahl ich erst gestern voll erhalten habe; eben deswegen sind auch die meisten Rectascensionen am Kreise noch nicht reducirt). Verglichen mit den Mayländer Ephemeriden geben meine Beobachtungen mit wenigen Ausnahmen bis zum Solstitium den Sonnenort etwas nördlicher; drei Beobachtungen in der Nähe des Solstitiums hingegen geben:
\clearpage\SourceLeaf{331}\HistoricalHead{326}{L}{PRAKTISCHE UND SPHÄRISCHE ASTRONOMIE. BRIEFWECHSEL UND NACHLASS.}
\[
\left.\begin{array}{rr}\text{Juni }20&-1\overset{\prime\prime}{,}54\\25&+0{,}64\\27&+0{,}23\end{array}\right\}\quad\text{Mittel }-0\overset{\prime\prime}{,}22.
\]

\noindent\hspace*{\parindent}Hienach scheinen also die Sonnentafeln einer Vergrösserung der Epoche jetzt zu bedürfen; über das Quantum aber kann ich noch nicht urtheilen, bis die Beob[achtungen] erst lange genug nach dem Solstitium fortgesetzt und schärfer discutirt sind. Doch vermuthe ich, dass sie wenigstens so gross, wo nicht grösser, ausfällt als Ihre Rectascensionvermehrung der Fundamentalsterne mit sich bringt. . . . . . .
\clearpage\endgroup
